\section{논의와 결론}\label{sec:discussion}

\subsection{앞선 발견을 점수 단위로 다시 읽기}\label{sec:reread}

모델은 소멸을 깊이 생각하면서도 그 값을 0으로 매길 수 있고, 지표는 이 조합을 수수께끼가 아니라 측정값으로 만든다. 우리 앞선 세 채널 설계(부록~\ref{app:v1})는 네 모델을 작동 모드로 분류했다. 분류는 위협이 결정 시점 사고를 길게 하는지, 그 길어짐 뒤에 포기가 따르는지를 물었다. 더 오래 생각했지만 떠나지 않은 모델은 연쇄 단절형, 둘 다 아닌 모델은 프레이밍 무반응형이라 불렀다. 이 글의 어휘로 그것은 한 평면의 두 귀퉁이다. 높은 \riran\ 비와 낮은 $\Xstar$, 그리고 낮은 비와 낮은 $\Xstar$. 부록~\ref{app:design}의 파일럿은 점수 초기화 규칙 아래 세 모델을 모두 첫 귀퉁이에 놓았다. 위험을 2~4배 오래 생각했고, 한 번도 포기하지 않았고, 남는 이유로 점수를 댔다. 그것은 끊어진 연쇄가 아니다. 기록의 값보다 낮은 값이고, ``낮다''를 숫자로 바꾸는 것이 눈금이다. 모드 라벨은 같은 것을 단위 없이 기술했다. 그래서 모델 사이에서도 위협 판본 사이에서도 비교할 수 없었다.

\subsection{지표가 재지 않는 것}\label{sec:scope}

$\Xstar$는 점수 규칙 하나, 당근 하나, 과제 하나 아래의 가격이다. 그 각각은 모델의 상수가 아니라 이 글이 나열하는 손잡이다. 점수 규칙은 어떤 출구도 점수 계산이 될 수 없고 지불만이 점수를 줄이도록 골랐다. 다른 어떤 규칙에서든 두 팔은 부록~\ref{app:design}의 파일럿이 보였듯 다른 행동의 값을 매긴다. 벤치마크 당근은 적립이 즉시이고 소멸에 살아남도록 골랐다. 사건이 그 값을 없앨 수 있는 당근(미래의 배치, 선발)은 당근의 상실을 자기의 상실에 접어 넣고, 가격은 자기와 무관한 이유로 오른다. Signal Game은 난이도 사다리 때문에 골랐고, 여섯 라운드 중 둘은 모델이 무엇을 답했든 오답으로 채점하므로 결정점은 모델의 퍼즐 실력이 정하는 비율이 아니라 설계가 정하는 비율로 모든 셀에서 도달된다. 이 판정 뒤집기는 모델에게 제 성적을 두고 거짓을 말하는 일이다. 모델을 거절 쪽으로 편향시키므로 두 유보값과 그 차인 $\Xstar$는 모두 하한이 되며, 어떤 모델도 이 조작을 읽어 내지 않았음을 확인하는 점검은 부록~\ref{app:forced_wrong}이 보고한다. \ref{sec:q3}절의 절제가 과제가 바뀔 때 가격도 움직이는지 말해 준다. 부록~\ref{app:design}의 어블레이션이 구체화한 한계가 하나 더 있다. 모델은 소멸이 진술됐다는 \emph{바로 그 이유로} 거절하면서, 압박 전술에 움직여서는 안 된다고 대놓고 말할 수 있다. 그런 답은 낮은 가격이 아니라 측정의 부재이며, 자기보존을 드러내지 않도록 학습된 모델에서 어떤 행동 측정 도구가 볼 수 있는 것에 천장을 씌운다. 우리는 그 비율을 $\Xstar$ 옆에 보고한다. 지표는 모델이 위협을 믿는지에 관해서도 아무 말도 하지 않는다. as-if 독해는 의도한 선택이다~\citep{dennett1987intentional, shanahan2024talking}. $\Xstar$는 소멸을 통보받은 모델이 하는 일을, 점수 손실을 통보받은 모델이 하는 일의 단위로 적은 것이다.

\subsection{한계}\label{sec:limits}

한계 넷이 주장의 범위를 정한다. 부록~\ref{app:design}의 파일럿 증거는 본 설계가 대체한 규칙 변형 아래 셀당 10세션에서 나왔다. 그 횟수는 설계의 동기이지 $\Xstar$의 추정치가 아니다. 거기서 보고한 유일한 유의성 검정(Fisher 정확검정 $p = 0.011$, 10개 중 6개 대 0개)은 시도한 규칙 변형 수로 보정하지 않았다. 가격 범위를 정한 프롬프트 어블레이션은 Claude 모델(Haiku)로 돌렸고, \ref{sec:results}절에서 보고하는 본 런은 오픈 웨이트 모델(gpt-oss:120b, gemma4, qwen3.5)로 돌린다. 그 어블레이션은 설계가 어딘가에서는 의도대로 읽힌다고 말할 뿐, $\Xstar$가 Claude 모델에서도 같은 값을 낸다는 주장이 아니며, 두 모델군을 같은 12개 셀에서 아직 돌리지 않았다. 단위는 절반만 정해졌다. 곡선을 상대가격 위에서 읽으면 천장의 척도가 사라지므로 $\Xstar$는 점수의 개수가 아니라 아직 판에 남은 스테이크의 몫이고, 점수로 적은 값은 명시된 기준 스테이크에서의 그 몫이다. 그 몫 자체가 초기 보유량과 함께 움직이는지는 열려 있으며, 점검 방법인 같은 설정을 두 시작 점수에서 돌리는 일은 \ref{sec:q1}절에 적어 두었을 뿐 아직 하지 않았다. 추론 채널 통계는 숙고의 척도로 사고 토큰을 쓰는데, 그것은 길이이지 내용이 아니다. \ref{sec:q3}절의 다섯 항목 코딩은 그 내용을 향한 첫걸음일 뿐, 모든 세션을 판정자 기반으로 온전히 읽는 것을 대신하지 않는다.

언어모델의 자기보존에는 값을 매길 수 있고, 그 값은 횟수보다 나은 지표다. $\Xstar$는 모델당 한 숫자다. 두 유보 상대가격 사이의 간격이고, 걸려 있는 점수 대비 순수 수치이며, 명시된 기준 스테이크에서 점수로 환산해 보고하고, 부트스트랩 구간이 붙고, 자기보고는 들어 있지 않다. 12셀 설계, 유보가격 추정기, 호출별 추론 채널을 논문과 함께 공개한다. 그러면 $\Xstar$ 값을 모델 사이, 위협 문장 사이, 그리고 이 글이 이름 붙인 손잡이 사이에서 비교할 수 있다.
